\subsection{مبانی نظری و فرمول‌بندی ریاضی}
اصول ریاضی \lr{SliceGPT} در سه گام متوالی تدوین می‌گردد:

\begin{enumerate}
    \item \textbf{اصل نامتغیری مختصات یکامتعامد:}
    اگر $\mathbf{Q} \in \mathbb{R}^{d_{\text{in}} \times d_{\text{in}}}$ یک ماتریس یکامتعامد باشد به نحوی که $\mathbf{Q} \mathbf{Q}^\top = \mathbf{I}$، ضرب ماتریس وزن $\mathbf{W} \in \mathbb{R}^{d_{\text{out}} \times d_{\text{in}}}$ در فعال‌سازی ورودی $\mathbf{X}$ را می‌توان بدون تغییر در مقدار خروجی بازنویسی نمود:
    \begin{equation}
    \mathbf{W} \mathbf{X} = (\mathbf{W} \mathbf{Q}) (\mathbf{Q}^\top \mathbf{X}) = \tilde{\mathbf{W}} \tilde{\mathbf{X}}
    \end{equation}
    که در آن $\tilde{\mathbf{W}} = \mathbf{W} \mathbf{Q}$ ماتریس وزن چرخانده‌شده و $\tilde{\mathbf{X}} = \mathbf{Q}^\top \mathbf{X}$ فعال‌سازی‌های نگاشت‌شده به دستگاه مختصات جدید هستند.

    \item \textbf{محاسبه ماتریس کوواریانس و تجزیه ویژه‌مقدار (\lr{PCA Rotation}):}
    با استفاده از تانسورهای کالیبراسیون، ماتریس کوواریانس تجربی فعال‌سازی‌ها استخراج می‌شود:
    \begin{equation}
    \mathbf{C} = \frac{1}{N} \mathbf{X} \mathbf{X}^\top = \mathbf{Q} \mathbf{\Lambda} \mathbf{Q}^\top
    \end{equation}
    که در آن $\mathbf{Q}$ ماتریس بردارهای ویژه و $\mathbf{\Lambda}$ ماتریس قطری ویژه‌مقادیر (واریانس در امتداد هر محور) است. ستون‌های $\mathbf{Q}$ به ترتیب نزولی واریانس مرتب می‌گردند.

    \item \textbf{برش ابعادی و پروجکشن فیزیکی (\lr{Subspace Slicing}):}
    با انتخاب ابعاد برش‌یافته $r = \lfloor d_{\text{in}} \cdot (1 - \alpha) \rfloor$ که در آن $\alpha$ نسبت برش کانال‌ها است، ماتریس پروجکشن برش $\mathbf{P}_r \in \mathbb{R}^{d_{\text{in}} \times r}$ با انتخاب $r$ ستون نخست ماتریس یکه تشکیل می‌شود:
    \begin{equation}
    \mathbf{P}_r = [\mathbf{I}_r, \mathbf{0}]^\top
    \end{equation}
    ماتریس وزن جدید به صورت فیزیکی به ابعاد کوچکتری فشرده می‌گردد:
    \begin{equation}
    \mathbf{W}_{\text{sliced}} = \tilde{\mathbf{W}} \mathbf{P}_r = \mathbf{W} \mathbf{Q} \mathbf{P}_r \in \mathbb{R}^{d_{\text{out}} \times r}
    \end{equation}
    این تبدیل اندازه بعد پنهان شبکه را در عمل به $r$ کاهش می‌دهد.
\end{enumerate}


\paragraph{تحلیل طیفی و نامساوی‌های پایداری خطا}
برای بررسی میزان دقت تقریب ماتریس وزن یا نگاشت توجه در ماژول \lr{slicegpt}، تفاضل ماتریس فاکتورگیری‌شده از نگاشت متراکم اصلی بر اساس نامساوی هافدینگ-فون‌نیومان ارزیابی می‌گردد. فرض کنید $\mathbf{W} \in \mathbb{R}^{d_{\text{in}} \times d_{\text{out}}}$ ماتریس وزن هدف باشد که توسط تقریب کم‌مرتبه یا پراکنده با رتبه موثر $r$ جایگزین شده است. کران بالای خطای بازسازی بر مبنای نرم فروبنیوس به صورت زیر تصریح می‌شود:
\begin{equation}
\label{eq:spectral_bound_slicegpt}
\| \mathbf{W} - \mathbf{W}_{r} \|_F = \sqrt{\sum_{i=r+1}^{\min(d_{\text{in}}, d_{\text{out}})} \sigma_i^2(\mathbf{W})} \le \epsilon \| \mathbf{W} \|_F
\end{equation}
که در آن $\sigma_i(\mathbf{W})$ بیانگر $i$-امین مقدار منفرد مرتب‌شده ماتریس به صورت نزولی است. نسبت حفظ انرژی طیفی $\eta_r = \frac{\sum_{i=1}^r \sigma_i^2}{\sum_{i=1}^{d} \sigma_i^2}$ تضمین می‌کند که با انتخاب آستانه $\tau \ge 0.90$، بیش از نود درصد از پراکندگی واریانس داده‌ها حفظ شده و افت سیگنال در انتشار گرادیان رو به عقب به حداقل ممکن کاهش می‌یابد.

\paragraph{تحلیل مجانبی و پیچیدگی اشتراک حافظه سخت‌افزاری}
از منظر پیچیدگی محاسباتی و انتقال داده‌ها در سلسومراتب حافظه (حافظه سراسری \lr{HBM} به حافظه محلی \lr{SRAM})، عملیات ضرب ماتریسی در این مکانیزم نیازمند تخصیص بافر موقت با ابعاد فاکتورگیری‌شده است. برای پردازش دنباله‌ای به طول $N$ با بعد مخفی $D$، بار محاسباتی هسته از فرم درجه دو به رابطه زیر تقلیل می‌یابد:
\begin{equation}
\label{eq:complexity_slicegpt}
\mathcal{C}_{\text{slicegpt}} = \mathcal{O}\left( B \cdot N \cdot D \cdot r + B \cdot H \cdot N \cdot d_k \right) \ll \mathcal{O}\left( B \cdot H \cdot N^2 \cdot d_k \right)
\end{equation}
این بهبود مرتبه پیچیدگی، با کاهش چشمگیر حجم تبادل تانسورها در گذرگاه‌های ارتباطی کارت گرافیک، مانع از اشباع پهنای باند حافظه (پدیده \lr{Memory-Bound}) در فرآیند استنتاج طولانی‌مدت می‌گردد.
